\section{第\ref{sec:recognition}章　認識}

認識の速さ、最初の段の仕組み、応答の線形な読み出し、時間の連続性による学習は、ヒトとサルで測定されている。連鎖全体を二つの演算に帰着させることと動画からの学習は本書のモデルで検証した。錯視の説明は、十分な根拠のあるものから論争中のものまである。

{\footnotesize
\setlength{\tabcolsep}{3pt}\renewcommand{\arraystretch}{1.15}
\begin{longtable}{>{\raggedright}p{0.5cm}>{\raggedright}p{4.0cm}>{\raggedright}p{5.6cm}>{\raggedright}p{2.3cm}>{\raggedright\arraybackslash}p{1.7cm}}
\toprule
No. & 主張 & 実験と測定されたもの & 出典 & 状態 \\
\midrule\endfirsthead
\toprule
No. & 主張 & 実験と測定されたもの & 出典 & 状態 \\
\midrule\endhead
1 & ずれがあると、画素ごとのテンプレート比較はコイン投げと変わらない；段の対~\eqref{eq:scstage}は図形をどの位置でも認識する & \texttt{models/\allowbreak recognition\_models.py}、$24$ 点の「網膜」、$4000$ 試行：反転した点の割合 $0$, $0.1$, $0.2$ でテンプレートは $0.49$, $0.51$, $0.52$；$S$ と $C$ の対は $1.00$, $0.86$, $0.70$ & 章内の導出 & モデル \\
2 & 絵の中の動物は $\sim150$\,ms で認識される。十数段を1段 $10$--$20$\,ms で1回通過する & ヒト、$20$\,ms 提示の写真で「動物がいるかいないか」の課題：二つの答えの誘発電位は $\sim150$\,ms で分かれる。サル：最初の応答の時刻は段を追って遅くなる（V1が最も早く、次にV2とV4）。段の数と「1段に0個か1個のスパイク」はこれらの数値からの推論 & \cite{Thorpe1996,Schmolesky1998} & 測定済み \\
3 & $S$ 段は部品の AND、$C$ 段は位置についての最大値（命題~\ref{prop:conveyor}） & ネコの一次視覚野：単純細胞は特定の場所の特定の向きのエッジに、複雑細胞は同じ向きにより広い領域で応答する。複雑細胞の細胞内記録：多くの細胞で2本のバーへの応答は二つの応答の大きいほうに近く、平均では max 演算に最も近い。相互抑制による最大値は命題~\ref{prop:wta}、全体の活動による除算は総説 & \cite{HubelWiesel1962,Lampl2004,Carandini2012} & 測定済み \\
4 & 連鎖を上るほど組み合わせは大きく複雑になり、応答は位置にますます依存しなくなる & サル、刺激を「臨界特徴」まで単純化：V4と後部下側頭皮質の細胞は小さな受容野で形、色、テクスチャの複雑な組み合わせに応答する；前部下側頭皮質では受容野が大きい。モデルの $S$ と $C$ の交互の繰り返しはこの描像を再現する & \cite{KobatakeTanaka1994,Riesenhuber1999} & 測定済み \\
5 & 畳み込みネットワークはこの交互の繰り返しを再現し、上位段の応答を最もよく予測する；物体はニューロン集団の発火率から線形に読み出せる & 脳へのあてはめなしに認識を学習したネットワーク：最上層はサル下側頭皮質のニューロン応答を、中間層はV4の応答を予測する。ランダムに選んだ $\sim100$ 個の細胞の $12.5$\,ms 以上の窓での活動から、線形分類器が位置や大きさが違っても物体とカテゴリーを認識する & \cite{Fukushima1980,Yamins2014,Hung2005} & 測定済み \\
6 & トレース付きヘッブ則~\eqref{eq:traceinvariance}は、トレースが $\sim20$\,ms 以上なら、ラベルなしの動画から不変性を学習する & 遅い特徴とトレース付き規則のアイデアは理論。\texttt{models/\allowbreak recognition\_models.py}、二つの $C$ ニューロン、入力 $16$、$10$ 回の実行、物体間の休止 $0.3$\,s。$\tau_{\text{tr}}=5$\,ms で図形との一致は平均 $0.52$、$10$\,ms で $0.56$；$20$, $50$, $200$\,ms ではすべての実行で $1.00$（$30$\,ms では10回中1回誤り） & \cite{Foldiak1991,WiskottSejnowski2002} & モデル \\
7 & 脳の中の不変性は時間の連続性から学習される & サル：視野の特定の場所への眼の跳躍中に物体をこっそり差し替えた；下側頭皮質ニューロンの位置不変性は差し替えに合わせて変わり、1時間後にはすでに有意だった。これが視覚学習の主要な規則だというのは仮説 & \cite{LiDiCarlo2008} & 測定済み \\
8 & 運動：方向検出器、次に広い領域についての同じ AND／OR の対 & ウサギ網膜の方向検出器；サルのMT野では記録した $110$ 個のニューロンすべてが方向選択的で、運動への応答はV1より強い & \cite{Barlow1965,Albright1984} & 測定済み \\
9 & 上位段は下へ予測を送り、上へは誤差が送られる & モデルは一次視覚野の受容野の古典外効果を説明する；個々の観察は整合する（総説）が、パイプラインの一般的な仕組みとしては証明されていない & \cite{RaoBallard1999,Keller2018} & 仮説 \\
10 & 難しい画像にはフィードバックと何周かの処理が要る & サル：「難しい」画像では下側頭皮質での決定が $\sim30$\,ms 遅れて形成される；この遅い応答は順方向ネットワークではうまく予測できず、フィードバックのあるネットワークやより深いネットワークのほうがよく予測する & \cite{Kar2019} & 測定済み \\
11 & 気づかれない画素の改ざんはネットワークをだます；ヒトの視覚ははるかに強いが無敵ではない & ネットワークで：特別に選んだ小さな摂動が答えを変える；「パンダ $\to$ テナガザル」の例は2番目の論文から。ヒトでは短時間提示のとき、ネットワーク間でよく転移する改ざんが選択をずらす & \cite{Szegedy2014,Goodfellow2015,Elsayed2018} & 測定済み \\
12 & 期待の錯視：信頼できない入力は期待に譲る。淡いものは「より遅く」動き、ドレスの見え方は人によって違う（\ref{sec:illusions}節） & コントラストの低い格子はより遅く動いて見える。$1401$ 人の調査：ドレスは三つの安定した見え方に分かれ、照明のヒントで色が切り替わる；$\sim13\,000$ 人の観察者では照明についての想定が決め手になる。遅い運動を期待するベイズモデルが一連の運動錯視を説明する & \cite{Thompson1982,LaferSousa2015,Wallisch2017,Weiss2002,Purves2011} & 仮説 \\
13 & ゲイン調整：滝、傾き、コントラスト；ヘルマン格子は近傍抑制ではない & ウサギ網膜の方向検出器は長い運動の間に発火率を下げ、停止後は背景以下に落ちる。傾きの錯視は近傍の活動による除算のモデルで再現される。網膜の出力ニューロンの応答を変えない格子の変形で斑点が消える。近傍抑制による説明は否定された & \cite{BarlowHill1963,Schwartz2009,SchillerCarvey2005} & 測定済み \\
14 & 遅延の補償：動く物体は閃光より前にあるように見える & 錯視は測定済み。心理物理は運動の前方延長とも遅延の差とも矛盾する；閃光後 $\sim80$\,ms の出来事にもとづく事後的な説明が提案されている。論争は決着していない & \cite{Nijhawan1994,EaglemanSejnowski2000} & 仮説 \\
15 & 静止した「ヘビ」が回る：遅延の差~\eqref{eq:latency}が方向検出器に読まれる & 錯視は色がなくても働く；隣り合う要素の対が単独で錯視を生む；方向選択的なサル皮質のニューロンは、この静止した対に方向性のある応答を示す。ヒトでは眼のマイクロサッカードと瞬きが回転を起動する & \cite{Conway2005,OteroMillan2012} & 測定済み \\
16 & 多義図形：相互抑制、疲労、ノイズ；切り替えは主にノイズが起こす & 競合するニューロン集団のモデル：切り替えを起こすのは\emph{ノイズ}であり、ノイズがなければ切り替えは起きない；刺激の強さとともに切り替え頻度が増すことを説明する。ここでは疲労の役割は小さい & \cite{MorenoBote2007} & 仮説 \\
17 & 錯視の一部はマシンが学習する課題から生じる & 動画の次フレームを予測するよう学習したネットワークは静止した輪の回転を「見て」、対照画像では見ない；画像のノイズ除去と鮮鋭化のネットワークはヒトと同様に色と明るさの錯視にかかる & \cite{Watanabe2018,GomezVilla2020} & 測定済み \\
\bottomrule
\end{longtable}}
